\appendix

% NARRATIVE: Negative results inform future work. These failures explain WHY
% WTA works (diversity in initialization, not architecture) by showing what
% doesn't work (diversity losses, slot attention, separate weights).
\section{Approaches That Failed}
\label{app:failed}

We document approaches that failed to improve over baseline, as negative results
inform future research directions. All experiments below achieved \emph{lower}
accuracy than standard TRM training.

\subsection{Diversity Through Training Losses}

Several approaches attempted to create diverse solution paths during training,
motivated by the success of halt-first ensembling at inference.

\paragraph{Forced trajectory divergence.} Adding a loss penalizing high cosine
similarity between trajectories from different input permutations. Despite
successfully reducing similarity (0.83 final), accuracy dropped from 86.3\% to
48.2\%---the model learned to satisfy both permutations with \emph{worse}
solutions rather than better diverse ones.

\paragraph{Inter-H diversity loss.} Penalizing consecutive H-iteration
similarity to prevent ``redundant'' iterations. This broke iterative
refinement: each step began undoing the previous step's work. Healthy models
show $\cos(z_H^{(h)}, z_H^{(h+1)}) > 0.9$; low similarity is harmful.

\paragraph{Multi-head output diversity.} Training $K{=}3$ separate value heads
on a shared trunk with diversity loss \emph{subtracting} similarity. Heads
collapsed to identical outputs within 1000 steps despite the diversity pressure.
The shared trunk creates a convergence attractor that overwhelms any loss signal.

\subsection{Architectural Diversity Mechanisms}

\paragraph{Slot attention.} Replacing $z_H$ with $K{=}3$ competing slots, forcing
cells to distribute information via softmax competition. Slots collapsed to
identical representations---the shared reasoning blocks (MLPMixer) dominate any
structural diversity mechanism.

\paragraph{K-specific full-rank weights.} Giving each of $K{=}3$ hypotheses
completely separate reasoning parameters (3$\times$ model size). Despite no
weight sharing, all hypotheses converged to the same solution. Separate weights
$\neq$ diversity; training dynamics must force divergence.

\paragraph{Learned $z_H$ initialization.} A puzzle-encoder MLP mapping inputs
to per-puzzle initial states. Accuracy dropped 1.9\%---fixed initialization
outperforms learned initialization, which adds noise and overfitting risk.

\paragraph{Multi-init training.} $K{=}3$ learned $(z_H, z_L)$ initialization pairs
with first-to-halt racing. Initial advantage vanished by step 2500; shared
weights caused all initializations to converge to the same solution.

\subsection{Loss Weighting and Curriculum}

\paragraph{Per-H label smoothing.} Softer targets early (exploration), sharper
late (commitment). Consistently 10+ percentage points behind baseline---high
average smoothing hurt convergence without compensating benefits.

\paragraph{Confidence-weighted loss.} Upweighting high-entropy cells. Initially
neutral but diverged late (--3\% at convergence); extra weighting destabilized
training.

\paragraph{Stuck-puzzle loss weighting.} 3$\times$ weight for puzzles where
$q_\text{halt} < 0.5$. With 91\% of puzzles marked ``stuck'' early in training,
this overwhelmed the loss landscape.

\subsection{Fixed-Point and Contraction Losses}

\paragraph{Fixed-point loss.} Penalizing $\|z_H^{(h+1)} - z_H^{(h)}\|$ to
encourage convergence. Accuracy dropped to ${\sim}65\%$. The model needs
continued iteration even when states are similar.

\paragraph{Validity loss.} Auxiliary loss encouraging predictions to satisfy
Sudoku constraints (no duplicates per row/col/box). All variants (symmetric,
low-temperature, Gumbel) failed at ${\sim}65\%$ accuracy.

\subsection{Input Manipulation}

\paragraph{Input corruption.} Corrupting 50\% of given digits during training
(diffusion-inspired). Destroyed learning---the model needs constraint information
intact. This contrasts with corrupting \emph{predictions} ($z_H$ noise), which
helps marginally.

\paragraph{Stochastic cell selection.} Updating random 50\% of cells per H-step
(ConsFormer-inspired). Failed by --3\%: TRM needs all cells to propagate
constraints simultaneously.

\subsection{Inference-Time Mechanisms at Training}

\paragraph{Binary reset during training.} Resetting $z_H/z_L$ when
$q_\text{halt} < 0.5$. DOA at step 500---destructive gradient signals from 41\%
of samples stuck at maximum iterations.

\paragraph{Entropy regularization.} Encouraging uncertainty at early H-steps.
Accuracy dropped 0.5\%: the model benefits from committing early on easy cells
and propagating, not from hedging.

\subsection{Architecture Modifications}

\paragraph{Causal attention.} Replacing MLPMixer with causal TransformerBlock.
Failed at --28\% behind baseline: TRM requires \emph{bidirectional} constraint
propagation. Causal masking prevents cell 80 from seeing cell 0, breaking
Sudoku's global structure.

\paragraph{Larger model.} Increasing hidden dimension from 512 to 768. Unstable
training with eventual collapse---the MLPMixer architecture doesn't benefit from
extra capacity on this task.

\subsection{Summary}

\begin{table}[h]
\centering
\small
\begin{tabular}{lc}
\toprule
Approach & $\Delta$ vs Baseline \\
\midrule
Forced trajectory divergence & --38\% \\
Causal attention & --28\% \\
Inter-H diversity loss & --20\% \\
Multi-head diversity loss & --15\% \\
Slot attention & --12\% \\
Per-H label smoothing & --11\% \\
Stochastic cell selection & --3\% \\
Confidence-weighted loss & --3\% \\
K-specific full-rank & --2\% \\
Learned $z_H$ init & --2\% \\
Entropy regularization & --0.5\% \\
\bottomrule
\end{tabular}
\caption{Summary of failed approaches, sorted by accuracy drop.}
\label{tab:failed}
\end{table}

The common failure mode: approaches that work for other architectures
(diffusion, ViTs, LLMs) often fail for TRM because iterative refinement has
different dynamics than single-pass inference or autoregressive generation.
TRM needs:
\begin{enumerate}
  \item Bidirectional information flow (no causal masking)
  \item High iteration-to-iteration similarity (incremental refinement)
  \item Early commitment on easy cells (propagation, not hedging)
  \item Intact input constraints (no corruption)
\end{enumerate}

The success of WTA training, in contrast, works \emph{with} these dynamics:
it creates diversity in the \emph{initialization} (where it matters) while
allowing each hypothesis to follow normal iterative refinement.
